\documentclass[11pt,a4paper]{article}

\usepackage[T1]{fontenc}
\usepackage[utf8]{inputenc}


\usepackage{amsmath,amssymb,amsthm}
\usepackage{mathtools}
\usepackage{hyperref}
\usepackage{doi}
\usepackage{geometry}
\usepackage{booktabs}
\usepackage{array}
\usepackage{enumitem}
\usepackage{xcolor}

\geometry{margin=2.5cm}

\hypersetup{
  colorlinks=true,
  linkcolor=black,
  citecolor=black,
  urlcolor=black,
  pdftitle={The Spectral Admissibility Sub-Programme --- Presentation Note 1},
  pdfauthor={J\'er\^ome Beau},
  pdfsubject={Map of the spectral admissibility sub-programme: constituent papers, phases,
    and the status of every result},
  pdfkeywords={spectral admissibility, Born--Infeld constraint, Weil representation,
    Heisenberg group, pair observable, capacity exponent, presentation note}
}

\theoremstyle{definition}
\newtheorem{remark}{Remark}
\newtheorem{definition}{Definition}

\newcommand{\Heis}{\mathrm{Heis}}
\newcommand{\SU}{\mathrm{SU}}
\newcommand{\su}{\mathfrak{su}}
\newcommand{\BFS}{\mathrm{BFS}}
\newcommand{\BI}{\mathrm{BI}}
\newcommand{\Cay}{\mathrm{Cay}}
\newcommand{\End}{\mathrm{End}}
\newcommand{\Sym}{\mathrm{Sym}}
\renewcommand{\Im}{\mathrm{Im}\,}
\newcommand{\dimpair}{\delta_{\mathrm{pair}}}
\newcommand{\betastar}{\beta^*}
\newcommand{\cchi}{c_{\mathrm{BI}}}
\newcommand{\nobs}{n^{\mathrm{obs}}_3}
\newcommand{\spair}{\sigma^{\mathrm{can}}_{\mathrm{pair}}}

\title{The Spectral Admissibility Sub-Programme\\[0.4em]
\large Presentation Note 1}

\author{J\'{e}r\^{o}me Beau\\[0.3em]
\small Independent Researcher\\
\small \texttt{jerome.beau@cosmochrony.org}}

\date{2026}

\begin{document}

  \maketitle

  \begin{abstract}
    The spectral admissibility sub-programme of the Cosmochrony corpus addresses a single
    central question: which spectral sectors of the Weil representation of
    $\Heis_3(\mathbb{Z}/q\mathbb{Z})$ survive the Born--Infeld bounded-flux constraint?
    The answer organises the entire O-series (papers O1 and O3--O33, five precursor papers, and
    the Span-Growth and Critical Coverage companion notes) along a single chain:
    \[
      \text{bounded flux}
      \;\Longrightarrow\;
      \text{admissibility envelope}
      \;\Longrightarrow\;
      \text{Weil-sector capacity}
      \;\Longrightarrow\;
      \dimpair .
    \]
    Its fibre-structure conditionality is closed by O24 for the segment
    $\cchi \to \dimpair$, under a supplied carrier.
    The further step to the cascade exponent $\betastar$ is not part of the chain: the
    reciprocal prescription $\betastar \approx 1/(\dimpair + \tfrac{1}{2})$ has no carrier on the
    Heisenberg measurement substrate, and the Span-Growth Note establishes that the
    expander-derived conversion does not transfer there, so the agreement of
    $\betastar \approx 0.126$ with the charged-lepton window is a coincidence check and not a
    derivation.
    On the SU(2) side the carrier is supplied conditionally rather than derived, and on the
    SU(3) side the O31 record is a withdrawal notice, so the colour group and the Standard Model
    gauge group are open.
    This note maps the constituent papers, identifies four internal phases, records the
    status of every result as proved, structural, numerical, conditional, or open, and states the
    open deliverables: hypothesis~$[\mathrm{H\text{-}color}]$, of which the two levels
    established in O32 stand while those that rested on O31 are not currently asserted; the
    identification of the spin-$\tfrac{1}{2}$ sector; and the extension of the numerical campaign
    to $q = 401$.
  \end{abstract}

  \medskip
  \noindent\textbf{Keywords.}
  spectral admissibility; Born--Infeld constraint; Weil representation; Heisenberg group;
  pair observable; capacity exponent; lepton mass hierarchy; SU(3); colour triplet;
  non-injective projection; presentation note.

  \newpage
  \tableofcontents

  \newpage

  \section{Motivation and Central Question}
    \label{sec:motivation}

    The spectral admissibility sub-programme answers the following question.

    \medskip
    \noindent\textit{Central question.} The non-injective projection $\Pi$ of the Cosmochrony
    framework acts on the Weil representation of $\Heis_3(\mathbb{Z}/q\mathbb{Z})$, decomposed
    into irreducible blocks $V_c$ indexed by characters $c \in (\mathbb{Z}/q\mathbb{Z})^\times$.
    The Born--Infeld saturation constraint bounds the projective flux carried by each mode:
    $A_n \leq A^{\max}_n := \cchi/\sqrt{\lambda_n}$, where $\lambda_n$ is the Laplacian
    eigenvalue at BFS depth $n$ and $\cchi$ is the BI saturation constant.
    \textit{Which combinations of Weil sectors remain admissible under this constraint, and
    what capacity exponent $\dimpair$ do they carry?}

    \medskip
    The practical interest of this question is that $\dimpair$ was proposed to govern the cascade
    exponent $\betastar$ through the reciprocal prescription
    $\betastar \approx 1/(\dimpair + \tfrac{1}{2})$.
    That prescription imports a changing-degree expander growth law into a fixed-degree Heisenberg
    cascade and has no native carrier there, and the Span-Growth Note~\cite{Beau2026sgn} proves the
    conversion does not transfer to the substrate on which the capacities are measured.
    The value $\betastar \approx 0.126$ matches the phenomenological window
    $\betastar \in (0.09, 0.13)$ derived from the charged-lepton mass ratios $m_e : m_\mu : m_\tau$;
    that agreement is a coincidence check rather than a quantitative connection between the
    projective framework and Standard Model observables.

    \medskip
    The sub-programme is also the hardest to navigate: it spans over thirty papers written
    over a period of successive obstruction identifications and resolutions.
    This note is intended as a structured entry point, not a summary of results.

  \section{Position in the Cosmochrony Programme}
    \label{sec:position}

    The Cosmochrony corpus is organised into three branches.
    Branch~I establishes the axiomatic primitive: admissible non-injective transitions as the
    primitive of physical description~\cite{Beau2026m}, with the structural necessity of
    non-injectivity~\cite{Beau2026l}.
    The selection of $\Heis_3(\mathbb{Z}/q\mathbb{Z})$ and its Weil representation from those
    axioms is \emph{not} a theorem: HeisenbergStructure~\cite{Beau2026HeisenbergStructure}
    disproves that implication by an explicit six-element countermodel, so the carrier is a
    supplied realisation on which this sub-programme's mathematics is conducted.
    Branch~II is the spectral admissibility sub-programme (this note).
    Branch~III derives quantum mechanics, spacetime geometry, gauge structure, and
    fermionic matter from Branch~I axioms and Branch~II spectral data.

    \medskip
    The spectral admissibility sub-programme therefore occupies a pivotal position: it is
    the \emph{computational engine} of the corpus.
    It takes the algebraic output of Branch~I (the Weil representation and its decomposition
    into irreducible blocks) and produces the two inputs that Branch~III depends on:
    \begin{enumerate}
      \item the identification of the admissible sector as the spin-$\tfrac{1}{2}$ representation
      of $\SU(2)$ embedded via the admissibility thread $Q_8 \subset 2I \subset \SU(2)$;
      \item the capacity exponent $\dimpair$ and its representation-theoretic interpretation as
      the scaling exponent of Hilbert--Schmidt norm growth along trajectories in the minimal
      admissible non-abelian sector $\su(2)$.
      The numerical value $\betastar \approx 0.126$ obtained from it by the reciprocal
      prescription is a coincidence check, not an output of the chain.
    \end{enumerate}

    \begin{remark}[Vocabulary note]
Throughout this note, the term \emph{admissibility constraint} replaces the earlier
``relaxation'' vocabulary of the preliminary Cosmochrony formulation.
The substrate $\chi$ is static: no substrate dynamics or relaxation mechanism is postulated
at the foundational level.
Observable evolution is an induced ordering of admissible projections, not a dynamics of
$\chi$ itself.
    \end{remark}

  \section{Logical Chain of the Sub-Programme}
    \label{sec:chain}

    The sub-programme is organised around the following derivation chain.

    \begin{equation}
      \label{eq:chain}
      \underbrace{\cchi}_{\text{BI saturation}}
      \;\Longrightarrow\;
      \underbrace{A^{\max}_n = \cchi/\sqrt{\lambda_n}}_{\text{admissibility envelope}}
      \;\Longrightarrow\;
      \underbrace{\spair(n)}_{\text{pair capacity}}
      \;\Longrightarrow\;
      \underbrace{\dimpair}_{\text{capacity exponent}}
      \;\;\text{\textemdash\textemdash}\;\;
      \underbrace{\betastar \approx 0.126}_{\text{prescription, no carrier}}.
    \end{equation}

    Each arrow in~\eqref{eq:chain} corresponds to a distinct analytical or numerical
    sub-problem, resolved at a different stage of the sub-programme.
    The final link is drawn without an arrow: it is a prescription, not a derivation.
    On the segment $\cchi \to \dimpair$, O24~\cite{Beau2026a28} closes the fibre-structure
    conditionality under the carrier supplied by O23, and no free parameter is adjusted.
    The step $\dimpair \to \betastar$ is not established: O24 states that it does not establish
    that capacity-to-rate step, for which no native Heisenberg growth carrier is defined; O25 and
    O14 record that the reciprocal prescription yields no $\betastar$ diagnostic from these data;
    and the Span-Growth Note~\cite{Beau2026sgn} proves that the expander-derived conversion law
    does not transfer to the Heisenberg substrate.
    The asymptotic numerical stabilisation of $\dimpair(q)$ as $q \to \infty$ is a
    separate analytical and numerical question, recorded as open in
    Section~\ref{sec:open-hyp}.
    The group-theoretic side carries conditionality on both sectors.
    The SU(2) identification rests on the carrier supplied conditionally by O23, and the SU(3)
    sector requires hypothesis~$[\mathrm{H\text{-}color}]$ while its O31 support is withdrawn
    (Section~\ref{sec:open}).

    \medskip
    Three structural facts underlie the chain and should not be confused.

    \begin{enumerate}
      \item The \emph{parity involution} $c \leftrightarrow q - c$, the candidate fibre
      structure of the pair sector (a structurally motivated hypothesis: O18~\cite{Beau2026a22}
      proves Born--Infeld parity as a covariance of the response family and states the fibre identification
      as an open classification problem), is the structural origin of the pair observable
      $\spair(n) = \sigma_c(n) \cdot \sigma_{q-c}(n)$.
      The factor-of-two relation $\dimpair = 2\delta_c$ is a consequence of this fibre pairing,
      conditional on the same hypothesis.

      \item The \emph{observable hierarchy}: the pair capacity $\sigma_c(n)$ (normalised BFS
      span in the Weil block $V_c$) is strictly finer than the Markov spectrum of the Weil
      operators.
      Spectral equality of characters does not imply capacity profile equality.
      This distinction is essential for hypothesis~$[\mathrm{H\text{-}color}]$.

      \item The \emph{admissibility thread} $Q_8 \subset 2I \subset \SU(2)$ is an output of
      the sub-programme, not an input.
      $Q_8$ is the minimal non-abelian finite subgroup of $\SU(2)$ compatible with bounded-flux
      admissibility; $2I$ is the maximal admissible finite subgroup; the LPS family provides the
      Ramanujan approximation to $\SU(2)$ in the large-$q$ limit.
    \end{enumerate}

  \section{Constituent Papers}
    \label{sec:papers}

    The sub-programme comprises five precursor papers, the O-series papers O1 and O3--O33, and two
    companion notes.
    Table~\ref{tab:papers} provides an overview organised by internal phase.

    \subsection{Precursor papers: SU(2) admissibility sub-programme}
      \label{sec:precursor}

      Four papers precede the O-series proper and explore spectral admissibility on Cayley graphs
      of finite $\SU(2)$ subgroups.

      \textbf{SpectralAdmissibility}~\cite{Beau2026a1} carries out the spectral admissibility
      programme on the chain $Q_8 \subset 2I \subset \mathrm{PSL}(2,\mathbb{F}_q)$.
      It derives the admissibility hierarchy $A^{\max}_n = \cchi/\sqrt{\lambda_n}$ in the
      linearised regime and shows the non-linear DBI correction preserves it.
      For the SU(3) chain $\Delta(27) \subset \Sigma(168)$, it establishes the co-admissibility
      of the fundamental and anti-fundamental representations via character theory.

      \textbf{SpectralCapacity}~\cite{Beau2026a2} introduces the capacity functional
      $\mathcal{C}_\alpha(G, S)$ and states the binary-polyhedral maximality \emph{conjecture}:
      that $2I$ uniquely maximises $\mathcal{C}_\alpha(G, S)$ for all $\alpha > 0$ among finite
      $\SU(2)$ subgroups with neutral generating sets.
      For the valences $d \in \{6, 12, 24\}$ the argument closes by a finite algebraic
      computation on classical character tables.
      The general statement remains a conjecture there: it would follow from a target theorem
      whose missing ingredient is a universal classification of the groups admitting a neutral
      generating set of a given size.

      \textbf{SpectralGramRigidity}~\cite{Beau2026a3} shows that the geometry of neutral
      generating sets on $\SU(2)$ is fully encoded in pairwise character values
      $G_{st} = -\tfrac{1}{2}\chi_\rho(st)$.

      \textbf{ThreeGenerations}~\cite{Beau2026a4} derives the three-generation structure from
      the ADE stratigraphic levels of binary Cayley graphs (structural).

      \textbf{SpectralRelaxation}~\cite{Beau2026a5} establishes $\Lambda_{\mathrm{proj}} \asymp
  \lambda_2$ for LPS expander families, providing the entry point for the O-series dynamics.

    \subsection{O-series core: LPS phase (O1--O8)}
      \label{sec:lps}

      The opening O-series papers work on Lubotzky--Phillips--Sarnak expander graphs, which have
      exponential shell growth $|B_n| \sim p^n$ and therefore represent the most natural
      large-group setting for bounded-flux dynamics.

      \textbf{O1}~\cite{Beau2026a6} shows that LPS saturation is a fixed-valence artefact;
      growing valence restores the pre-saturation window.

      \textbf{O3}~\cite{Beau2026a7} derives the phenomenological target window
      $\betastar \in (0.09, 0.13)$ from the charged-lepton mass ratios $m_e : m_\mu : m_\tau$
      (structural).

      \textbf{O4}~\cite{Beau2026a8} examines whether a structural upper bound $\beta \leq 1$
      follows from the BI flux constraint together with the Cheeger isoperimetric inequality, and
      finds that it does not follow from that argument.

      \textbf{O5}~\cite{Beau2026a9} tests vertex-based class-function cascades on
      $\mathrm{PSL}(2,\mathbb{F}_q)$ and finds that none of the constructions it examines yields a
      viable mechanism for $\betastar$.

      \textbf{O6}~\cite{Beau2026a10} confirms $q$-structural saturation at the $O(q^2)$
      Steinberg level and characterises the obstruction to exponent extraction there: the
      pre-saturation window is too short to extract a stable effective exponent, and a multi-step
      path fingerprint is identified as the construction the next level would require.

      \textbf{O7}~\cite{Beau2026a11} reformulates the target as a capacity exponent
      $\delta \in [7.4, 10.6]$ on three-step path fingerprints (structural).

      \textbf{O8}~\cite{Beau2026a12} establishes the geometric obstruction closing the LPS
      phase: exponential shell growth compresses the entire pre-saturation window into
      $O(\log q)$ BFS steps, preventing stable slope estimation.

    \subsection{O-series core: Heisenberg transition (O9--O15)}
      \label{sec:heisenberg}

      The geometric obstruction of O8 forces a change of arena.
      The Cayley graph $\Cay(\Heis_3(\mathbb{Z}/q\mathbb{Z}), S_q)$ has polynomial ball growth
      $|B_n| \sim C n^4$ (Bass--Guivarc'h theorem, homogeneous dimension $D = 4$), restoring an
      exploitable window whose depth O9 bounds below by $\Omega(q^{1/2})$ BFS steps.

      \textbf{O9}~\cite{Beau2026a13} resolves the geometric obstruction by switching to
      Heisenberg graphs.
      Polynomial growth is confirmed, and its Window-Depth Theorem shows that a vertex window of
      $O(q^2)$ spans $\Omega(q^{1/2})$ BFS steps; the diameter is $\Theta(q)$.

      \textbf{O10}~\cite{Beau2026a14} performs the first large-$q$ computation ($q = 101$,
      $|G_q| = 1{,}030{,}301$) and identifies an algorithmic obstruction: the dense TensorSketch
      fingerprint is structurally mismatched to the problem, requiring sketch dimension
      $D \gg q^4$.

      \textbf{O11}~\cite{Beau2026a15} resolves the algorithmic obstruction via a bidimensional
      Fourier-character proxy for the Weil-block decomposition; a stable estimate
      $\hat\delta_{\mathrm{cap}} \approx 3.4$ is extracted at $q = 307$.

      \textbf{O12}~\cite{Beau2026a16} implements the exact Weil projection, closes the
      representation-level ambiguity left by O11, and extracts $\hat\delta_{\mathrm{exact}}
  \approx 4.3$--$4.8$ for $q \in \{29, 53, 61\}$ ($R^2 > 0.98$).

      \textbf{O13}~\cite{Beau2026a17} confirms asymptotic stability of the exact Weil-block
      capacity over the extended prime range and establishes monotone decrease of
      $\hat\delta_{\mathrm{exact}}(q)$, definitively refuting the finite-size hypothesis.

      \textbf{O14}~\cite{Beau2026a18} establishes why exact Heisenberg capacity does not
      determine the phenomenological cascade rate.
      A factor $q^{-\eta}$ changes the intercept of a fixed-$q$ regression and leaves its fitted
      slope invariant, the proposed central-phase correction vanishes identically because the
      central coordinate multiplies each exact Fourier mode by a scalar phase and is algebraically
      invisible to rank, and the reciprocal prescription yields no $\betastar$ diagnostic from
      these data.

      \textbf{O15}~\cite{Beau2026a19} bounds the weighted aggregation observable: for its
      dynamical weights, $\alpha_{\mathrm{dyn}} \leq \hat\delta_{\mathrm{exact}}$ whenever a
      checkable covariance condition holds over the fitting window (Corollary~4.6 of that paper),
      verified there at $q = 151$ and $q = 211$ and open at the other tested primes.
      Since $\hat\delta_{\mathrm{exact}} < 5.0$ at every tested prime, that observable cannot
      reach $[7.4, 10.6]$ where the condition is confirmed.
      O15 states that the stronger claim, that no non-negative reweighting whatsoever could exceed
      $\hat\delta_{\mathrm{exact}}$, is false, and that its restricted result does not rule out
      every conceivable reweighting.
      It is what motivates the pair observable in O16.

    \subsection{Pair observable and transfer chain (O16--O24)}
      \label{sec:pair}

      O15 also establishes the Heisenberg transfer no-go of its title: any future resolution
      requires a new growth carrier or an explicitly stated cross-substrate hypothesis.
      Its aggregation bound points, in addition, at the observable class.
      The physically correct observable is the \emph{pair capacity}
      $\spair(n) = \sigma_c(n) \cdot \sigma_{q-c}(n)$ on conjugate pairs $\{c, q-c\}$.

      \textbf{O16}~\cite{Beau2026a20} examines whether conjugate blocks are the fibres of $\Pi$
      and types the two premises that identification would need, neither derived here.
      Under a further hypothesis, that the conjugate-pair ratio $r(c,q)$ is independent of the
      shell index, the log-log slope of the pair product is exactly double that of $\sigma_c$.
      On the exact block observable of the O12/O13 pipeline the proved conjugate-block identity
      forces $r(c,q) \equiv 1$, so the pair product reduces definitionally to the single block
      squared and there is no separate doubling phenomenon there.
      O16 asserts no numerical value of the pair exponent.

      \textbf{O17}~\cite{Beau2026a21} proves the conjugation identity $\rho_{q-c} = \bar\rho_c$
      and the exact block-independence of the Gram--Schmidt count, within its own scalar toy
      model; the reading of the conjugate pair as a fibre remains a hypothesis.

      \textbf{O18}~\cite{Beau2026a22} delimits what Born--Infeld structure contributes to that
      hypothesis: the action is even and parity acts covariantly on the response family
      (even-order responses preserved, odd-order reversed; proved), an even
      countermodel refutes fixed-probe indiscernibility, and the fibre identification is stated
      as a typed open classification problem (O18 Problem~2.8), with the Weil-level bridge
      $c \leftrightarrow q-c$ open.

      \textbf{O19}~\cite{Beau2026a23} proves that the residual amplitude factor
      $r(c, q) \in \{1, 2, 3, 4\}$ does not affect $\dimpair$; the observable is
      pipeline-independent.

      \textbf{O20}~\cite{Beau2026a24} establishes a dynamical selection criterion for admissible
      modes via projective persistence (structural).

      \textbf{O21}~\cite{Beau2026a25} defines the observable saturation rank $\nobs$ via the
      threshold condition $\Sigma_c(n_3) = 3$, states a conditional fibre-to-rate prescription
      rather than a constructed transfer, and reads the O15 bound as an observable-class
      mismatch.

      \textbf{O22}~\cite{Beau2026a26} proves that Born--Infeld projection locking forces
      $\nobs$ onto a BFS shell of $\Cay(G_q, S_q)$: shell-alignment is a theorem, not an
      independent geometric hypothesis.

      \textbf{O23}~\cite{Beau2026a27} determines the status of the complementary question: for a
      supplied irreducible two-dimensional unitary carrier, the traceless neutral sector is
      $\mathfrak{su}(2) \cong \Im\mathbb{H}$, of real dimension exactly three (Theorem~3.1).
      The carrier selection and the identification of $\Sigma_c$ with that dimension are open
      bridges, so the threshold value $\Sigma_c(n_3) = 3$ is a supplied selection rule, not a
      derived constant.

      \textbf{O24}~\cite{Beau2026a28} is the capstone of the analytical O-series.
      The Verticality Lemma (``$\ker\Pi$ is free, $\Im\Pi$ is constrained'') closes the
      fibre-structure conditionality of the segment $\cchi \to \dimpair$: its Corollary~4.2 states
      that the observable-rank condition supporting that segment is independent of the cardinality
      of the fibres of $\Pi$, conditional on the carrier supplied by O23 and with the fibre
      identification itself open.
      O24 states explicitly that this does not establish the separate $\dimpair \to \betastar$
      capacity-to-rate step, for which no native Heisenberg growth carrier is defined.

    \subsection{Numerical campaign and closure papers (O25--O33), with the companion notes}
      \label{sec:numerical}

      With the fibre-structure conditionality of $\cchi \to \dimpair$ closed by O24, the
      post-O24 programme turns to three tasks:
      numerical validation of $\dimpair$, representation-theoretic identification of the
      admissible sector, and the SU(3) colour sector.

      \textbf{O25}~\cite{Beau2026a29} carries out the systematic pair-level numerical campaign
      across all $(q-1)/2$ conjugate pairs with $M = 50$ independent block samples per pair, for
      $q \in \{29, 61, 101, 151, 211\}$.
      Results: $\bar\delta_{\mathrm{pair}} \in \{8.89, 9.98, 9.55, 9.13, 8.78\}$ with inter-pair
      standard deviation decreasing from 0.54 to 0.11.
      The campaign is extended there to $q \in \{307, 401, 601\}$ at reduced sampling.
      After the finite-size adjustment $\eta \log q / \log n_1(q)$, the adjusted values
      $\delta_{\mathrm{corr}}(q) \in [7.7, 9.0]$ lie within the target window $[7.4, 10.6]$ for
      $q \in \{61, 101, 151, 211\}$, and O25 reports the adjustment overcorrecting beyond that
      range, to $6.53$ at $q = 601$.
      O25's own reading assigns the window agreement to the \emph{raw} exponent, which descends
      into $[7.4, 10.6]$ unaided and reaches $7.61$ at $q = 601$.
      O14~\cite{Beau2026a18} establishes that this term leaves a fixed-$q$ fitted slope invariant,
      so the adjustment is a declared diagnostic and not a correction of the fitted exponent.

      \textbf{O26}~\cite{Beau2026a30} establishes the quadratic completion of admissible
      spectral pairs via binary-icosahedral representation, providing the Level~III
      identification criterion for the spin-$\tfrac{1}{2}$ sector.

      \textbf{O27}~\cite{Beau2026a31} shows that every admissible morphism factors uniquely
      through the admissibility quotient, and exhibits the canonical representation action
      $\rho \colon \su(2) \to \End(V_\rho)$ under two premises: the identification of the
      admissible image with the neutral traceless sector, which O27 assumes and derives nowhere,
      and the spinor carrier supplied conditionally by O23.
      The intertwiner space is one-dimensional, so an equivariant map is unique only up to a
      scalar, and no non-zero equivariant map reaches the fundamental from the adjoint, so the
      vector lift of O26 Hypothesis~4.4 remains open.
      The identification chain reads, with arrows of unequal type:
      \[
        \text{emergence} \Rightarrow \su(2)\;[\text{O27, conditional}]
        \Rightarrow r_{\mathrm{eff}} = 3\;\text{[O28, measured; O29, audited]}.
      \]
      No step of it identifies a representation carrier.

      \textbf{O28}~\cite{Beau2026a32} calibrates the BFS window out of sample, where the fitted
      linear extrapolation fails, and reports the adjusted $\delta_{\mathrm{corr}}(q)$ within
      $[7.4, 10.6]$ over $q \in \{29, 61, 101, 151, 211\}$.
      In the samples of its perimeter the covariance $C_c$ in $\End(H_{\mathrm{eff}})$ has $1\%$
      threshold rank $3$ for every conjugate pair at $q \in \{61, 101, 151\}$, with resolved
      spectrum $[1 : \tfrac{1}{2} : \tfrac{1}{2}]$.
      That value is finite-data and not an invariant: the extended recomputation of O29 finds rank
      three modal, at $32/50$ pairs for $q = 101$ and $103/105$ for $q = 211$.

      \textbf{O29}~\cite{Beau2026a33} audits what that measurement identifies, and reaches no
      representation.
      It separates the measured trajectory span in $H_{\mathrm{eff}}$, a data-dependent leading
      truncation inside it, and the spin-$\tfrac{1}{2}$ module $V_\rho$, which is supplied
      elsewhere.
      Its finite-precision diagnostics are compatible with an adjoint reading of
      $H_{\mathrm{eff}}$ without deriving that carrier or proving exact irreducibility, and no
      Veronese square root is recovered there.
      The doublet dimension $d_\rho = 2$ is the minimality selection of O26, not a measurement of
      these data.

      \textbf{O30}~\cite{Beau2026a34} gives a conditional $3\times3$ model on $\Sym^2(V_\rho)$ in
      which the ratio $\lambda_1/\lambda_2 = 2$ follows from an equatorial constraint, under its
      hypotheses (H0)--(H3) and, for the spectrum, two weighted decorrelation conditions.
      The factor of two is the squared symmetric-tensor normalisation of the direction the Carnot
      grading identifies as central.
      That model is not the $9\times9$ covariance measured in O28, and O30 supplies no analytical
      account of it.

      \textbf{O31}~\cite{Beau2026a35} is a withdrawal notice.
      It withdraws the uniqueness of $\SU(3)$ among connected compact subgroups acting irreducibly
      on $\mathbb{C}^3$, the irreducibility of the group of Born--Infeld-compatible
      fibre-preserving transformations, the co-admissibility of the individual profiles, the
      arithmetic criterion tying colour triplets to $q \equiv 1 \pmod 3$, the metaplectic identity
      as it was used, the $\Delta(27)$ spectral proposition, a mistyped rank-eight test in
      $\End(\mathbb{C}^3)$, and every derivation of a gauge factor or of the Standard Model gauge
      group as a composite.
      Three counterexamples are given there; $\mathrm{SO}(3) \subset \SU(3)$ alone defeats the uniqueness
      claim.
      The scalar, arithmetic and discrete-group material of the superseded version, including the
      generator-twist isospectrality statement, is under re-audit there and is not asserted.
      The notice records that any dependent claim elsewhere resting on a withdrawn result is to be
      treated as unsupported until separately re-examined.

      \textbf{O32}~\cite{Beau2026a36} provides the numerical test of $[\mathrm{H\text{-}color}]$
      on the standard graph.
      Its measurements stand as measurements; they do not support the withdrawn colour-group
      derivation.
      Results for $q \in \{61, 151, 211\}$ (test, $q \equiv 1 \pmod 3$) and
      $q \in \{29, 101\}$ (control, $q \equiv 2 \pmod 3$): hypothesis~$[\mathrm{H\text{-}color}]$
      passes all three test primes; $q = 307$ confirms ($R_{\mathrm{var}} = 0.7$, all 4 triplets).

      \textbf{O33}~\cite{Beau2026a37} gives a finite spectral architecture for mass-square
      splitting and mixing: exact doublets and protected sectors from Weil admissibility.

      \subsubsection*{Companion notes}

      \textbf{Span-Growth Note}~\cite{Beau2026sgn} classifies the native span growth induced by
      capacity decay and proves that the expander-derived capacity-to-rate conversion does not
      transfer to the Heisenberg measurement substrate, for the corpus-defined closure, frontier,
      pair-observable and exponent coordinates.
      It is the result that types the last link of~\eqref{eq:chain} as a prescription without a
      carrier.

      \textbf{Critical Coverage Note}~\cite{Beau2026ccnote} resolves the O25--O28 saturation-depth
      asymptotics by an exact interval theorem: $n_1(q) \to 22$ in probability and
      $x_1(q) = |B_{n_1}|/q^2 \asymp q^{-2}$, with deterministic bounds for every prime
      $q \geq 311$.

      \begin{table}[ht]
        \centering
        \caption{Summary of the O-series by internal phase.
        Status codes: P = proved; S = structural; N = numerical; C = conditional; O = open.}
        \label{tab:papers}
        \renewcommand{\arraystretch}{1.2}
        \begin{tabular}{@{}>{\raggedright\arraybackslash}p{2.7cm}
          >{\raggedright\arraybackslash}p{3.1cm}
          >{\raggedright\arraybackslash}p{7.3cm}l@{}}
          \toprule
          Phase & Papers & Central output & Status \\
          \midrule
          Precursors & SpAdm, SpCap, SpGram, 3Gen & $Q_8 \subset 2I$; binary maximality for
            $d \in \{6, 12, 24\}$ & P/S \\
          LPS phase & SpRel, O1--O8 & Geometric obstruction; no cascade-exponent bound & P/O \\
          Heisenberg transition & O9--O15 & Exact $\hat\delta_{\mathrm{exact}} \approx 4.5$;
            conditional aggregation bound & P/C \\
          Pair + transfer & O16--O24 & $\cchi \to \dimpair$ fibre-conditionality closed & C \\
          Numerical + sector & O25--O30 & $\delta_{\mathrm{corr}}$ diagnostic;
            $r_{\mathrm{eff}} = 3$ finite-data & N/C \\
          SU(3) / colour & O31--O32 & O31 withdrawn; O32 measurements stand & O \\
          Mass architecture & O33 & Finite spectral architecture for mass-square splitting & S \\
          Companion notes & Span-Growth, Critical Coverage & Transfer failure; exact depth law & P \\
          \bottomrule
        \end{tabular}
      \end{table}

  \section{Inputs from Upstream Results}
    \label{sec:inputs}

    The spectral admissibility sub-programme takes the following inputs from Branch~I.

    \begin{enumerate}
      \item \textbf{Weil representation.}
      $\Heis_3(\mathbb{Z}/q\mathbb{Z})$ and its Weil representation, taken as a supplied carrier.
      The two-dimensional module $V_\rho$ of the later sections is a further supplied object,
      selected by O26 minimality, and is not the Weil representation itself.
      HeisenbergStructure~\cite{Beau2026HeisenbergStructure} disproves their selection from the
      axioms A1--A4 by an explicit six-element countermodel, so the carrier is a supplied
      realisation and the mathematics of this sub-programme is conducted on it.

      \item \textbf{BI saturation functional.}
      The Born--Infeld action with saturation constant $\cchi$ is the unique analytic completion
      compatible with bounded projective transport capacity.
      The sub-programme takes $\cchi$ as a parameter; its structural origin is in Branch~I.

      \item \textbf{Non-injectivity no-go.}
      The ENI paper~\cite{Beau2026l} establishes that non-injectivity is a structural necessity
      of genuine emergence, not an interpretive option.
      The pair structure and the candidate parity involution $c \leftrightarrow q - c$ are
      downstream constructions, the fibre identification remaining an open problem.

      \item \textbf{Noscale principle.}
      The absence of an external dimensional scale at the admissibility level~\cite{Beau2026j}
      ensures that the only dimensionful parameter available to the sub-programme is $\cchi$,
      enforcing the internal consistency of the capacity law.
    \end{enumerate}

    \noindent No result from Branch~III is an input to the sub-programme.
    The admissibility thread $Q_8 \subset 2I \subset \SU(2)$ and the value $\dimpair$ are
    outputs; $\betastar$ is not, the conversion to it having no carrier on this substrate.

  \section{Outputs to Downstream Branches}
    \label{sec:outputs}

    Table~\ref{tab:outputs} lists the principal outputs of the sub-programme and their
    downstream consumers.

    \begin{table}[ht]
      \centering
      \caption{Principal outputs of the spectral admissibility sub-programme and their consumers.
        Status codes as in Table~\ref{tab:papers}.}
      \label{tab:outputs}
      \renewcommand{\arraystretch}{1.3}
      \begin{tabular}{@{}>{\raggedright\arraybackslash}p{5.5cm}
        >{\raggedright\arraybackslash}p{4.5cm}l@{}}
        \toprule
        Output & Consumer & Status \\
        \midrule
        $\dimpair$; the reciprocal $\betastar \approx 0.126$ is a coincidence check
          & SMchargemass, lepton hierarchy & N \\
        Admissibility thread $Q_8 \subset 2I \subset \SU(2)$ & Q1--Q3 (quantum sector) & P \\
        Spin-$\tfrac{1}{2}$ sector $V_\rho \cong \mathbb{C}^2$ (supplied, not derived) & Q7--Q11 (geometry) & C \\
        $\Sigma_c(n_3) = 3$ and $\Im\mathbb{H} \cong \su(2)$ (supplied carrier) & Q7 (spatial identification) & C \\
        SU(3) co-admissible triplet & Q6a, Gauge note & withdrawn \\
        $\delta_{\mathrm{corr}}(q) \in [7.4, 10.6]$ for $q \in \{61,\ldots,211\}$,
          a declared adjustment & Open: $q \to \infty$ limit & N \\
        \bottomrule
      \end{tabular}
    \end{table}

  \section{Status of Results}
    \label{sec:status}

    \subsection{Proved results}
      \label{sec:proved}

      The following results are established as theorems within their stated assumptions.

      \begin{enumerate}
        \item \textit{Binary-polyhedral maximality at the verified valences}
        (SpectralCapacity~\cite{Beau2026a2}): for $d \in \{6, 12, 24\}$, $2I$ maximises
        $\mathcal{C}_\alpha(G, S)$ for all $\alpha > 0$, by a finite algebraic computation.
        The statement for general $d$ remains a conjecture there.

        \item \textit{Gram rigidity} (SpectralGramRigidity~\cite{Beau2026a3}): the geometry of
        neutral generating sets on $\SU(2)$ is fully encoded in pairwise character values.

        \item \textit{Aggregation bound} (O15~\cite{Beau2026a19}): for its dynamical weights and
        under a checkable covariance condition, $\alpha_{\mathrm{dyn}} \leq
        \hat\delta_{\mathrm{exact}}$, so that observable cannot reach $[7.4, 10.6]$ where the
        condition holds.
        The universal reweighting no-go is not established.

        \item \textit{Parity covariance} (O18~\cite{Beau2026a22}): BI action evenness and
        the $(-1)^{k}$ covariance of the response family are proved; fixed-probe
        indiscernibility is refuted, and the
        minimal fibre condition $\Pi^{-1}(x) = \{\chi, -\chi\}$ is an open classification
        problem, not a forced consequence.

        \item \textit{Pipeline independence of $\dimpair$} (O19~\cite{Beau2026a23}): the
        residual amplitude factor $r(c,q)$ does not affect the pair exponent.

        \item \textit{Shell-alignment theorem} (O22~\cite{Beau2026a26}): Born--Infeld projection
        locking forces $\nobs \in \mathbb{N}$ as the first element of $L_{\mathrm{BI}} \cap N_q$.

        \item \textit{Conditional adjoint dimension} (O23~\cite{Beau2026a27}): for a supplied
        irreducible $\mathbb{C}^2$ carrier, $\dim_\mathbb{R} \Im\mathbb{H} = 3$ (Theorem~3.1);
        the threshold $\Sigma_c(n_3) = 3$ is a supplied selection rule, its carrier-selection
        and $\Sigma_c$-identification bridges being open.

        \item \textit{Observable-rank stability} (O24~\cite{Beau2026a28}): the Verticality Lemma
        closes the fibre-structure conditionality of the segment $\cchi \to \dimpair$, the
        observable-rank condition being independent of fibre cardinality.
        It remains conditional on the carrier supplied by O23, and it does not establish the
        step $\dimpair \to \betastar$.

        \item \textit{Admissible factorisation} (O27~\cite{Beau2026a31}): every admissible
        morphism factors uniquely through the admissibility quotient.
        Identifying that quotient with $\su(2)$ is conditional on the neutral-sector identification
        O27 assumes and on O23's supplied carrier, so it does not belong to this list.
      \end{enumerate}

    \subsection{Structural results}
      \label{sec:structural}

      The following results are derived from structural arguments that are internally consistent
      but not yet fully formalised as theorems.

      \begin{enumerate}
        \item \textit{Phenomenological window} (O3~\cite{Beau2026a7}): $\betastar \in (0.09, 0.13)$
        from $m_e : m_\mu : m_\tau$.

        \item \textit{Capacity exponent target} (O7~\cite{Beau2026a11}): the target
        $\delta \in [7.4, 10.6]$ on three-step path fingerprints.

        \item \textit{Observable-class identification} (O16~\cite{Beau2026a20}): the pair
        observable resolves the gap by exponent doubling $\dimpair = 2\delta_c$.

      \end{enumerate}

      No SU(3) emergence framework is listed: it would rest on the metaplectic identity that
      O31~\cite{Beau2026a35} withdraws.

    \subsection{Conditional results}
      \label{sec:conditional}

      The following hold under premises supplied elsewhere and not derived within the
      sub-programme.

      \begin{enumerate}
        \item \textit{The $\su(2)$ identification} (O27~\cite{Beau2026a31}, O23~\cite{Beau2026a27}):
        conditional on the identification of the admissible image with the neutral traceless
        sector and on O23's supplied spinor carrier.

        \item \textit{The doublet dimension $d_\rho = 2$} (O26~\cite{Beau2026a30}): a minimality
        selection, not a measurement of the checkpoint data.

        \item \textit{The eigenvalue ratio $\lambda_1/\lambda_2 = 2$} (O30~\cite{Beau2026a34}):
        in a conditional $3\times3$ model on $\Sym^2(V_\rho)$, under (H0)--(H3) and two weighted
        decorrelation conditions.
        It is not an account of the $9\times9$ covariance measured in O28.

        \item \textit{The adjoint reading of $H_{\mathrm{eff}}$} (O29~\cite{Beau2026a33}):
        supplied rather than derived; the finite-precision diagnostics are compatible with it and
        do not establish it.
      \end{enumerate}

    \subsection{Numerical evidence}
      \label{sec:numerical-evidence}

      The following results are numerically established but not yet derived analytically.

      \begin{enumerate}
        \item \textit{$\delta_{\mathrm{corr}}(q) \in [7.7, 9.0]$} (O25~\cite{Beau2026a29}): for
        $q \in \{61, 101, 151, 211\}$ after the finite-size adjustment; within the target window
        $[7.4, 10.6]$.

        \item \textit{Effective rank} (O28~\cite{Beau2026a32}, O29~\cite{Beau2026a33}):
        $r_{\mathrm{eff}} = 3$ for every conjugate pair at $q \in \{61, 101, 151\}$ in O28's
        perimeter; modal rather than invariant in O29's extended recomputation, at $32/50$ pairs
        for $q = 101$ and $103/105$ for $q = 211$.

        \item \textit{Eigenvalue structure $[1:\tfrac{1}{2}:\tfrac{1}{2}]$} (O28~\cite{Beau2026a32}):
        resolved across the measured pairs, with no analytical derivation.

        \item \textit{$[\mathrm{H\text{-}color}]$ numerical support} (O32~\cite{Beau2026a36}):
        passes for $q \in \{61, 151, 211, 307\}$ with variance ratio $R_{\mathrm{var}} \leq 1$
        relative to the control noise floor.
      \end{enumerate}

    \subsection{Open hypotheses}
      \label{sec:open-hyp}

      The following items remain open within the sub-programme.

      \begin{enumerate}
        \item \textit{$[\mathrm{H\text{-}color}]_{\mathrm{pointwise}}$}: open.
        The single-frequency argument for it belonged to the superseded version of O31 and is not
        asserted by the current record~\cite{Beau2026a35}, which is a withdrawal notice; nothing
        there refutes the argument, and no result of the sub-programme currently establishes it.

        \item \textit{Asymptotic stabilisation of $\delta_\infty$}: the limit
        $\delta_\infty = \lim_{q \to \infty} \dimpair(q)$ has not been computed analytically.
        The correct asymptotic variable is $n_1(q)/q$; its law is settled by the Critical Coverage
        Note~\cite{Beau2026ccnote}, which gives $n_1(q) \to 22$ in probability and
        $x_1(q) \asymp q^{-2}$, so $n_1(q)/q \to 0$.

        \item \textit{Extension to $q = 401$}: the campaign at $q \in \{29, 61, 101, 151, 211, 307\}$
        is complete; $q = 401$ is pending.
      \end{enumerate}

  \section{Open Deliverables}
    \label{sec:open}

    Three open deliverables define the boundary of the sub-programme as of this note.

    \textbf{Open deliverable 1: $[\mathrm{H\text{-}color}]$.}
    Of its four levels, two stand and two do not.
    (ii)~$[\mathrm{H\text{-}color}]_{\mathrm{avg}}$ is proved in O32 (Proposition~1) to
    $O(q^{-1})$, and (iii)~$[\mathrm{H\text{-}color}]_{\mathrm{eff}}$ in O32 in the
    $q \to \infty$ limit; both stand.
    (i)~$[\mathrm{H\text{-}color}]_{\mathrm{rank}}$ and
    (iv)~$[\mathrm{H\text{-}color}]_{\mathrm{pointwise}}$ rested on the superseded version of
    O31~\cite{Beau2026a35}, whose current record asserts neither; they are open.
    No $\SU(3)$ identification follows on either graph.

    \textbf{Open deliverable 2: the spin-$\tfrac{1}{2}$ identification.}
    The carrier is supplied by O23 and selected by O26 minimality; the vector lift of O26
    Hypothesis~4.4 is open, and no measurement of the sub-programme identifies $V_\rho$.

    \textbf{Open deliverable 3: the numerical campaign at full breadth.}
    The pair-level campaign runs at $q \in \{29, 61, 101, 151, 211\}$ at full breadth and is
    extended in O25 to $q \in \{307, 401, 601\}$ at reduced sampling ($M = 16$ at $307$, $M = 8$
    and twelve pairs at $401$ and $601$), so the large-$q$ dispersion there is inflated by the
    reduced pair count.
    At $q = 307$ the O32 variance test passes on all 4/4 triplets
    ($R_{\mathrm{var}} = 1.17 \times 10^{-3}$,~\cite{Beau2026a36}).
    What remains is breadth rather than reach: a full-breadth run beyond $q = 211$ would
    constrain the convergence of $n_1(q)/q$ and supply an uninflated data point for the
    $R_{\mathrm{var}}(q) \propto q^{-1}$ scaling.

    \medskip
    \noindent All other open problems of earlier papers (e.g., the rigorous continuum limit of
    the state law, the exact form of $\Phi$ on LPS graphs) are sub-problems within already
    resolved phases and do not block the downstream branches.

  \section{Conclusion}
    \label{sec:conclusion}

    The spectral admissibility sub-programme provides the core computational content of the
    Cosmochrony corpus.
    Starting from the Born--Infeld bounded-flux constraint on the Weil representation of
    $\Heis_3(\mathbb{Z}/q\mathbb{Z})$, taken as a supplied carrier, it measures the capacity
    exponent $\dimpair$, selects the spin-$\tfrac{1}{2}$ admissible sector by minimality, and
    establishes the admissibility thread $Q_8 \subset 2I \subset \SU(2)$ as a structural output.

    The SU(2) sector is carried by a conditional chain rather than closed: its carrier is supplied
    by O23 and the vector lift of O26 Hypothesis~4.4 is open.
    The numerical campaign is complete for $q \in \{29, 61, 101, 151, 211, 307\}$.
    The extension to SU(3) is open: O31~\cite{Beau2026a35} withdraws the colour-group derivation
    and every derivation of a gauge factor, so neither the colour factor nor the Standard Model
    gauge group is established here; the O32 measurements~\cite{Beau2026a36} stand on their own
    terms.

    The outputs of this sub-programme are used throughout Branch~III: by the quantum sector
    (Q1--Q3), by the geometric emergence sub-programme (Q5a--Q11), by the gauge structure
    sub-programme (Q6a, O31--O32), and by the fermionic matter papers (Q14); consumers that
    relied on the withdrawn O31 results are to treat them as unsupported until re-examined.
    Reading the O-series without this map is possible but significantly more difficult;
    this note is intended to make the logical structure transparent before the technical
    machinery is engaged.

    \newpage
    \bibliographystyle{plain}
    \bibliography{cosmochrony-bibliography}

\end{document}
